% Chapter 7: 逐企深度拆解
\section{L0 核心生态枢纽}

\subsection{协创数据（300857）--- \starfive{}}

\begin{table}[H]
\centering\small
\renewcommand{\arraystretch}{1.3}
\begin{tabularx}{\textwidth}{l X}
\toprule
\textbf{维度} & \textbf{详情} \\
\midrule
与英伟达关系 & A股唯一同时拥有 Jetson + DRIVE双平台 + NCP云服务商 三重顶级授权 \\
确定性证据 & Omnibot平台原生接入Cosmos 3全栈；独家承接Coalition国内落地部署 \\
最新业绩 & 2026Q1营收60.85亿(+192.9\%)，净利7.5亿(+343.45\%) \\
投资逻辑 & Cosmos 3国内落地的"总代理+方案商"，卡位生态核心入口，弹性最大 \\
\bottomrule
\end{tabularx}
\end{table}

\subsection{中科创达（300496）--- \starthree{}$\frac{1}{2}$}

\begin{table}[H]
\centering\small
\renewcommand{\arraystretch}{1.3}
\begin{tabularx}{\textwidth}{l X}
\toprule
\textbf{维度} & \textbf{详情} \\
\midrule
与英伟达关系 & 长期软件生态伙伴，OS定制/边缘AI \\
风险 & \textcolor{riskred}{2025.3美子公司列入BIS实体清单}；实控人质押71\% \\
投资逻辑 & 系统软件工程服务费+授权，非概念型但需关注北美客户影响 \\
\bottomrule
\end{tabularx}
\end{table}

\section{L1 算力底座}

\subsection{工业富联（601138）--- \starthree{}$\frac{1}{2}$}

\begin{table}[H]
\centering\small
\renewcommand{\arraystretch}{1.3}
\begin{tabularx}{\textwidth}{l X}
\toprule
\textbf{维度} & \textbf{详情} \\
\midrule
与英伟达关系 & 全球一级核心代工商，GB200/GB300机柜主力生产商 \\
最新业绩 & 2026Q1营收2511亿(+56.52\%)，净利106亿(+102.55\%) \\
重大风险 & \textcolor{riskred}{2024.10 BIS实体清单（24子公司）+TDO拒绝令覆盖223家子公司} \\
投资逻辑 & 算力"卖水人"，业绩确定性最强，但BIS风险需控仓位 \\
\bottomrule
\end{tabularx}
\end{table}

\section{L2 物理仿真与合成数据}

\subsection{索辰科技（688507）--- \starfive{}}

\begin{table}[H]
\centering\small
\renewcommand{\arraystretch}{1.3}
\begin{tabularx}{\textwidth}{l X}
\toprule
\textbf{维度} & \textbf{详情} \\
\midrule
与Cosmos 3关系 & "天工·开物"可微分物理AI平台接入Cosmos 3推理框架 \\
确定性证据 & 2025年报104次提及物理AI；物理AI业务已创收5816万元 \\
风险提示 & Q1亏损3394万；8月解禁11.7\% \\
投资逻辑 & A股最接近"物理规律软件化"的标的 \\
\bottomrule
\end{tabularx}
\end{table}

\subsection{五一世界（06651.HK）--- \starfive{}}

\begin{table}[H]
\centering\small
\renewcommand{\arraystretch}{1.3}
\begin{tabularx}{\textwidth}{l X}
\toprule
\textbf{维度} & \textbf{详情} \\
\midrule
与英伟达关系 & Omniverse生态首批中国合作伙伴，智驾仿真龙头 \\
主体澄清 & \textcolor{riskred}{NuRec（纽瑞克斯）是另一家独立公司，非51WORLD旗下} \\
投资逻辑 & "物理AI第一股"标签，港股流动性弱但弹性大 \\
\bottomrule
\end{tabularx}
\end{table}

\section{L3 感知硬件}

\subsection{奥比中光（688322）--- \starfive{}}

\begin{table}[H]
\centering\small
\renewcommand{\arraystretch}{1.3}
\begin{tabularx}{\textwidth}{l X}
\toprule
\textbf{维度} & \textbf{详情} \\
\midrule
与英伟达关系 & Jetson平台官方3D视觉合作伙伴 \\
确定性证据 & Gemini 330完成Cosmos 3数据闭环适配；CES 2026联合展出 \\
最新业绩 & Q1扣非净利+531\%，机器人收入达上年同期3倍以上 \\
投资逻辑 & 物理AI感知入口---机器人的"眼睛" \\
\bottomrule
\end{tabularx}
\end{table}

\section{L4 执行终端}

\subsection{德赛西威（002920）--- \starfive{}}

\begin{table}[H]
\centering\small
\renewcommand{\arraystretch}{1.3}
\begin{tabularx}{\textwidth}{l X}
\toprule
\textbf{维度} & \textbf{详情} \\
\midrule
与英伟达关系 & 国内智驾核心战略客户，DRIVE Thor芯片首发权+首家中国Tier1 \\
确定性证据 & IPU14高阶智驾域控2026Q1量产；双Thor+NVLink 4000TFLOPS \\
订单数据 & 域控市占率33.6\%（国内第一）；在手智驾订单超130亿元 \\
风险提示 & Q1营收-4.37\%，净利-20.74\%，短期承压 \\
投资逻辑 & 智驾域控绝对龙头，Thor量产后业绩弹性最大（240亿增量） \\
\bottomrule
\end{tabularx}
\end{table}

\subsection{绿的谐波（688017）--- \starfour{}}

\begin{table}[H]
\centering\small
\renewcommand{\arraystretch}{1.3}
\begin{tabularx}{\textwidth}{l X}
\toprule
\textbf{维度} & \textbf{详情} \\
\midrule
与物理AI关系 & Cosmos 3缩短研发→加速量产→减速器需求爆发（间接） \\
重大澄清 & \textcolor{riskred}{公司2025.6/12月多次公告否认与Tesla Optimus有任何直接业务关系} \\
最新业绩 & Q1营收1.40亿(+42.96\%)，净利3263万(+61.17\%)（已核验） \\
投资逻辑 & 国产谐波替代+人形机器人减速器需求；与Optimus叙事需区分 \\
\bottomrule
\end{tabularx}
\end{table}

\subsection{天准科技（688003）--- \starfive{}}

\begin{table}[H]
\centering\small
\renewcommand{\arraystretch}{1.3}
\begin{tabularx}{\textwidth}{l X}
\toprule
\textbf{维度} & \textbf{详情} \\
\midrule
与英伟达关系 & Jetson平台国内精英合作伙伴 \\
确定性证据 & 基于Jetson Thor具身大脑域控制器（2026.2发布）已大批量销售 \\
最新业绩 & Q1营收3.84亿(+75.61\%)，扭亏为盈 \\
投资逻辑 & 物理AI硬件执行层最直接受益者，产品最成熟 \\
\bottomrule
\end{tabularx}
\end{table}

\section{L2 物理仿真（续）}

\subsection{智微智能（001339）--- \starfour{}}

\begin{table}[H]
\centering\small
\renewcommand{\arraystretch}{1.3}
\begin{tabularx}{\textwidth}{l X}
\toprule
\textbf{维度} & \textbf{详情} \\
\midrule
与英伟达关系 & Jetson NPN合作伙伴（有授权）；基于Isaac SIM开发机器人虚拟仿真系统 \\
确定性证据 & "智擎"系列机器人大脑域控制器已量产；智算业务毛利率84.71\%（纯软件） \\
最新业绩 & Q1营收+53.29\%，净利+158.76\% \\
风险提示 & \textcolor{riskred}{9月首发限售解禁6.2\%，5家创投排队减持；机构可投性D级} \\
投资逻辑 & 把Isaac SIM/Cosmos本土化+硬件化的边缘算力角色 \\
\bottomrule
\end{tabularx}
\end{table}

\section{L3 感知硬件（续）}

\subsection{柯力传感（603662）--- \starfour{}}

\begin{table}[H]
\centering\small
\renewcommand{\arraystretch}{1.3}
\begin{tabularx}{\textwidth}{l X}
\toprule
\textbf{维度} & \textbf{详情} \\
\midrule
与物理AI关系 & 力控传感→物理交互数据闭环底层；与Cosmos 3无直接合作，属刚需硬件 \\
确定性证据 & 六维力送样80+家客户，月出货破千只；上证e互动官方确认已进入特斯拉供应链+有订单落地 \\
供应链范围 & 已落地：①电池包视觉检测②上海工厂轴重秤③三维力小批量；送样中：车载端六维力（高价值） \\
体量说明 & \textcolor{riskamber}{特斯拉相关营收占比$<$3\%，主要为工厂端检测设备；车载端大规模定点存在不确定性} \\
投资逻辑 & 7月Optimus投产是关键催化；月出货若达3000只→年化4.5亿利润 \\
\bottomrule
\end{tabularx}
\end{table}

\section{L4 执行终端（续）}

\subsection{中大力德（002896）--- \starthree{}$\frac{1}{2}$}

\begin{table}[H]
\centering\small
\renewcommand{\arraystretch}{1.3}
\begin{tabularx}{\textwidth}{l X}
\toprule
\textbf{维度} & \textbf{详情} \\
\midrule
关联路径 & 英伟达 → 宇树H2+ GR00T参考设计 → 宇树供应链（第一大行星减速器供应商） \\
确定性证据 & 32.6亿框架订单（2025正式1.8亿+2026意向8亿+2027意向22.8亿）；宇树批量供货互动易确认 \\
份额说明 & \textcolor{riskamber}{63\%份额公开信披未查到（华经情报网约25\%），宇树IPO 2026.1终止无法验证} \\
\textcolor{riskred}{重大合规风险} & \textcolor{riskred}{2025.3证监会处罚1500万+实控人10年禁入：虚增营收5.59亿/利润1.57亿+大股东占资4.62亿；2024.9立案/2025.4 ST；实控人质押85.4\%} \\
风险提示 & Q1净利-36.42\%；30.8亿为意向协议非正式合同；叠加信披造假前科，订单可信度需打折 \\
投资逻辑 & 赌宇树IPO催化；\textcolor{riskred}{ESG红牌，禁止入库（HRF-8触发）} \\
\bottomrule
\end{tabularx}
\end{table}

\subsection{鸣志电器（603728）--- \starthree{}$\frac{1}{2}$}

\begin{table}[H]
\centering\small
\renewcommand{\arraystretch}{1.3}
\begin{tabularx}{\textwidth}{l X}
\toprule
\textbf{维度} & \textbf{详情} \\
\midrule
关联路径 & 英伟达 → 宇树量产 → 空心杯电机/步进电机需求 \\
确定性证据 & 宇树科技空心杯电机第一大供应商，已获大额订单 \\
最新业绩 & Q1营收+18.05\%，净利+92.03\% \\
\textcolor{riskred}{BIS制裁} & \textcolor{riskred}{2024.6美子公司AMP被BIS列入实体清单（占总资产/营收$\approx$20\%），2025.3拟分拆AMP申请移除} \\
\textcolor{riskred}{监管警示} & \textcolor{riskred}{2024.7上海证监局警示函：业绩预告差-54\%未修正+关联方占用1.36亿；12月上交所通报批评} \\
投资逻辑 & 空心杯独家但PE 300$\times$纯情绪票；与英伟达无直接合作；需跟踪AMP分拆BIS移除进展 \\
\bottomrule
\end{tabularx}
\end{table}

\subsection{天娱数科（002354）--- 2星（高危）}

\begin{table}[H]
\centering\small
\renewcommand{\arraystretch}{1.3}
\begin{tabularx}{\textwidth}{l X}
\toprule
\textbf{维度} & \textbf{详情} \\
\midrule
与Cosmos 3关系 & \textcolor{riskred}{无直接英伟达官方合作/授权，仅买GPU自用} \\
\textcolor{riskred}{重大合规风险} & \textcolor{riskred}{2024.4因营收不足1亿+亏损被ST；2024.10因信披违规被证监会立案（第二次）；2024.12实控人徐德伟被山东监委留置；2025.1收到处罚事先告知书（虚增营收嫌疑）；2025.10已摘帽} \\
投资逻辑 & \textcolor{riskred}{强烈不建议参与}。追涨风险极高。ESG禁止入库（HRF-9触发） \\
\bottomrule
\end{tabularx}
\end{table}
